\appendix
\chapter{General Coding Standards \& Guidelines}
\label{app:coding-standards}

To ensure the maintainability, readability, and scalability of the codebase for Medicare -- A Smart Healthcare Reminder App, the following coding standards and guidelines were followed throughout the development using Flutter and Firebase.

\section{Variable Naming Conventions}
Consistent naming is used throughout the app:
\begin{itemize}
    \item \texttt{snake\_case}: Used for file names and folder names (e.g., \texttt{medication\_screen.dart}).
    \item \texttt{PascalCase}: Used for class and widget names (e.g., \texttt{MedicationReminder}, \texttt{HomeScreen}).
    \item \texttt{camelCase}: Used for variables and function names (e.g., \texttt{reminderTime}, \texttt{saveReminder()}).
\end{itemize}

\section{Visual Scope Indicators}
Variables are named clearly to show their scope:
\begin{itemize}
    \item \texttt{\_privateVariable} (underscore means it is private in Dart)
    \item \texttt{const MAX\_RETRIES = 3}
    \item \texttt{final String userId}
\end{itemize}

\section{Read-Only and Immutable Fields}
\begin{itemize}
    \item The \texttt{final} keyword is used for values that do not change after being set.
    \item The \texttt{const} keyword is used for values that are fixed at compile time.
    \item Example: \texttt{final String userId = auth.currentUser!.uid}
\end{itemize}

\section{Function Naming and Behavior}
All function names clearly describe what they do:
\begin{itemize}
    \item \texttt{saveMedicationReminder()}
    \item \texttt{loadUserData()}
    \item \texttt{scheduleNotification()}
    \item \texttt{markIntakeAsTaken()}
    \item \texttt{fetchPatientRecords()}
\end{itemize}

\section{Commenting Style}
Comments are added where the code is not easy to understand.
\begin{lstlisting}
// This function saves the medication reminder to Firestore
Future<void> saveMedicationReminder(Reminder reminder) async {
    // Validate fields before saving
}
\end{lstlisting}

\section{Indentation and Code Structure}
\begin{itemize}
    \item 2-space indentation is used as per Flutter/Dart standard.
    \item Each widget and function is clearly separated for easy reading.
    \item Related code is grouped inside the same file or class.
\end{itemize}

\section{File and Folder Naming}
All file names use \texttt{snake\_case} as required by Dart:
\begin{itemize}
    \item \texttt{medication\_screen.dart}
    \item \texttt{notification\_helper.dart}
    \item \texttt{chatbot\_screen.dart}
\end{itemize}
Folders are organized by feature: \texttt{auth/}, \texttt{medication/}, \texttt{measurement/}, \texttt{activity/}, \texttt{appointment/}, \texttt{chatbot/}, \texttt{admin/}.

\section{Widget and Class Structure}
\begin{itemize}
    \item Each widget or class handles only one responsibility.
    \item UI widgets are separated from business logic.
    \item Provider or setState is used for state management.
    \item Firebase calls are handled inside repository classes, not inside UI widgets.
\end{itemize}

\section{Access Modifiers and Scoping}
\begin{itemize}
    \item Private variables use underscore prefix (e.g., \texttt{\_reminderTime}).
    \item Public methods are only made public when needed by another class.
    \item Local variables are declared inside functions to limit their scope.
\end{itemize}

\section{Cohesion and Coupling}
\begin{itemize}
    \item Each module (Medication, Measurement, Activity, etc.) works independently.
    \item Shared utilities like \texttt{NotificationHelper} and \texttt{FirebaseHelper} are used across modules without tight dependency.
\end{itemize}

\section{DRY Principle (Don't Repeat Yourself)}
\begin{itemize}
    \item Common logic is written once and reused across modules.
    \item Helper classes like \texttt{FirebaseHelper} and \texttt{InputValidator} are shared across the whole app.
    \item Repeated UI components are made into reusable widgets.
\end{itemize}

\section{Single Responsibility Principle (SRP)}
\begin{itemize}
    \item Each part of the code has one job.
    \item A function handles a single task.
    \item A class handles a single responsibility (e.g., \texttt{ReminderScheduler}, \texttt{UserAuth}).
    \item A screen file only handles the UI for that screen.
\end{itemize}

\section{Extensibility and Maintainability}
\begin{itemize}
    \item New modules can be added without changing existing code.
    \item Shared base classes and utility functions make future updates easy.
    \item Firebase structure is designed to support adding new features easily.
\end{itemize}

\section{Exception Handling and Logging}
\begin{itemize}
    \item All Firebase and network calls are wrapped in try-catch blocks.
    \item Errors are shown to the user with simple error messages using SnackBar.
\end{itemize}

\begin{lstlisting}
try {
    await firestore.collection('reminders').add(data);
} catch (e) {
    print('Failed to save reminder: $e');
}
\end{lstlisting}
